% Complete readable proofs of the checked manuscript arguments.
% Reviewed source main.tex SHA256:
% af9cf13785178b750bb359975ca0f649961191ddd52a20d05a599a0ce2bdf508
% This is an appendix fragment, using the manuscript's existing notation.
% These are proofs of the stated results following the Lean development;
% proof-method correspondence is stated explicitly for each assertion.

\subsection{One-tile bound by the written H\"older argument (ANT-4-1)}
\label{leanproof:ANT-4-1}

The proof below establishes both inequalities of Lemma~\ref{lem:tile-height}
with its exact coefficient, by the multiple-function H\"older inequality and
translation invariance used in the manuscript. Lean also checks all
measurability, integrability, normalization, and height-coordinate changes.

\begin{proof}[Complete proof following the Lean argument]
Fix \(r\ge1\), \(c<1/2\), \(\varepsilon>0\), and a tile
\(B=\varepsilon\Delta'\). Write \(a=1/2-c>0\). The function
\(w=\psi_E e^\sigma\) is positive and continuous, and \(B\) is compact.
Consequently \(w_B=\sup_Bw\) is finite, and
\(\prod_{k=1}^r w(\mathbf f^{(k)})\le w_B^r\) on \(B^r\).
The integrand is nonnegative, so this proves the first inequality once
the integrals are known to be finite. We prove that finiteness together
with the second inequality.
The floor regions and height fibers are measurable: strict coordinate
comparability is open, so its exclusion is closed, and the coordinate
maps are continuous. All kernels used below are nonnegative measurable
functions.

For fixed floors \(\mathbf f=(\mathbf f^{(1)},\ldots,\mathbf f^{(r)})\),
use the common height \(s\) and relative heights
\[
 \boldsymbol\delta=(0,\delta_2,\ldots,\delta_r),\qquad
 h^{(1)}=s,\quad h^{(k)}=s+\delta_k\quad (k\ge2).
\]
The map \((s,\delta_2,\ldots,\delta_r)\mapsto(h^{(1)},\ldots,h^{(r)})\)
has determinant \(1\). Adding the same height to every point changes
no pairwise difference. Thus the allowed relative heights are exactly
the pinned fiber \(\mathcal H_r^0(\mathbf f)\), independently of \(s\).
For \(r=1\) the relative-height space is zero dimensional and has
its usual unit volume.

Put \(g_a(u)=\exp(-au^2)\). Equal-exponent H\"older for the \(r\)
nonnegative measurable functions \(s\mapsto g_a(s+\delta_k)\) gives
\[
 \int_{\mathbb R}\prod_{k=1}^r g_a(s+\delta_k)\,ds
 \le \prod_{k=1}^r\left(\int_{\mathbb R}g_a(s+\delta_k)^r\,ds\right)^{1/r}
 =\int_{\mathbb R}e^{-ar s^2}\,ds
 =\sqrt{\frac\pi{ar}}.
\]
Here translation invariance makes every factor identical; the \(r\)th roots
multiply back to the same integral, including \(r=1\).
For full precision the formal proof first applies H\"older to nonnegative
extended-real integrals, so it needs no prior integrability assumption.
The finite Gaussian value on the right then proves integrability of the
product. The elementary exponential product identity and normalization yield
\[
 (2\pi)^{-r/2}\int_{\mathbb R}
     \exp\!\left(-a\sum_{k=1}^r(s+\delta_k)^2\right)\,ds
 \le (2\pi)^{-r/2}\sqrt{\frac\pi{ar}}
 =\frac{(2\pi)^{-(r-1)/2}}{\sqrt{r(1-2c)}}.
\]
This is the manuscript's argument with the constant
\((2\pi)^{-1/2}\) factored out of its function
\(\phi(u)e^{cu^2}=(2\pi)^{-1/2}g_a(u)\).

Tonelli's theorem and the determinant-one change of height variables
now yield
\[
 \begin{aligned}
 &(2\pi)^{-r/2}\int_{\mathcal A_r(B)}
                  e^{-a\|\mathbf h\|^2}\,d\mathbf f\,d\mathbf h\\
 &\quad=
   \int_{B^r}\int_{\mathcal H_r^0(\mathbf f)}
       \left((2\pi)^{-r/2}\int_{\mathbb R}
          e^{-a\sum_k(s+\delta_k)^2}\,ds\right)
       d\boldsymbol\delta\,d\mathbf f\\
 &\quad\le
   \frac{(2\pi)^{-(r-1)/2}}{\sqrt{r(1-2c)}}
   \int_{B^r}\operatorname{Vol}_{r-1}
                       (\mathcal H_r^0(\mathbf f))\,d\mathbf f\\
 &\quad=
   \frac{(2\pi)^{-(r-1)/2}}{\sqrt{r(1-2c)}}\,V_r(B).
 \end{aligned}
\]
The previously proved pinned-volume bound makes the last expression
finite. Hence the middle integral is an ordinary integrable
nonnegative function, rather than an undefined integral assigned a
default value.

The affine coordinate-permutation symmetries of the actual reflection
tiles preserve pinned volume. Scaling \(r\) floor variables and \(r-1\)
free heights has Jacobian
\(\varepsilon^{(d-1)r+r-1}=\varepsilon^{dr-1}\). Therefore
\[
 V_r(B)=V_r(\varepsilon\Delta')
       =\varepsilon^{dr-1}V_r(\Delta).
\]
Substitution proves that the middle integral is at most
\(v_r(\varepsilon,c)\). Multiplying by \(w_B^r\) proves the second
inequality and, at the same time, integrability of the weighted
integrand by domination. The separately defined empty coefficient
\(v_0=1\) is consistent with the empty configuration; the displayed
common-height argument applies to \(r\ge1\).
\end{proof}

\begingroup\small\raggedright
The exact Lean steps are
\nolinkurl{Antichain.lintegral_translated_product_holder},
\nolinkurl{Antichain.lintegral_translated_gaussian_holder},
\nolinkurl{Antichain.translated_gaussian_integrable_holder},
\nolinkurl{Antichain.integral_translated_gaussian_product_le}, and
\nolinkurl{Antichain.normalized_common_height_lintegral_le}
in \nolinkurl{CommonHeightHolder.lean} and
\nolinkurl{OneTileGaussian.lean};
\nolinkurl{Antichain.heightCoordinates_measurePreserving} and
\nolinkurl{Antichain.antichainHeightFiber_lintegral}
in \nolinkurl{HeightCoordinates.lean};
and \nolinkurl{Antichain.antichainHeightFiber_gaussian_lintegral_le}
in \nolinkurl{HeightGaussianBound.lean}.
The actual tile, scaling, and weight-supremum estimates are combined in
\nolinkurl{Antichain.oneTile_model_full_spec}.
The finite ordinary-integral endpoints are
\nolinkurl{Antichain.oneTile_real_full_spec},
\nolinkurl{Antichain.oneTile_middle_integrableOn}, and
\nolinkurl{Antichain.oneTile_weighted_integrableOn}
in \nolinkurl{OneTileReal.lean}.
\par\endgroup

\subsection{Poissonization and the written spatial-independence product bound (ANT-4-3)}
\label{leanproof:ANT-4-3}

The proof below follows the manuscript's finite-intensity Poisson process,
independent restrictions to disjoint tiles, and finite-union exhaustion.
The spatial independence is proved from the actual sampled count and marks:
it is not an assumed property or a replacement by an allocation inequality.
The allocation calculation appears inside its proof and accounts for the
factorials explicitly.

\begin{proof}[Complete proof following the Lean argument]
Fix \(0<\varepsilon\le1\), \(c<1/2\), and bounded continuous
\(\sigma:E\to\mathbb R\). Put \(a=1/2-c>0\) and
\(w=\psi_Ee^\sigma\). On \(X=E\times\mathbb R\), let
\[
 d\mu(\mathbf f,h)
    =w(\mathbf f)\frac{e^{-ah^2}}{\sqrt{2\pi}}\,d\mathbf f\,dh.
\]
A mark \((\mathbf f,h)\) represents the point
\(\mathbf f+h\mathbf e\) in \(\mathbb R^d\).
The orthogonal floor-height change of variables has unit Jacobian;
thus \(\mu\) is exactly the weighted Gaussian point measure used to
define \(\mathcal Z_n(\sigma,c)\). Gaussian integration gives
\[
 0<\Lambda:=\mu(X)
    =\frac{1}{\sqrt{1-2c}}\int_Ew(\mathbf f)\,d\mathbf f<\infty.
\]
The finiteness follows from boundedness of \(\sigma\) and
integrability of \(\psi_E\).

For any finite measure \(\nu\) on \(X\), define
\[
 A_n(\nu)=
   \int_{X^n}\mathbf 1_{\{\text{the represented points form an antichain}\}}
                  \,d\nu^{\otimes n},\qquad A_0(\nu)=1,
\]
and
\[
 G_\nu(t)=\sum_{n\ge0}\frac{t^nA_n(\nu)}{n!}\qquad(t\ge0).
\]
These are nonnegative integrals and series. Moreover,
\(A_n(\nu)\le\nu(X)^n\), so \(G_\nu(t)\le e^{t\nu(X)}<\infty\).
The antichain events are measurable because they exclude finitely
many open strict-comparability conditions.
By the actual product-density formula,
\[
 A_n(\mu)=\mathcal Z_n(\sigma,c).
\]

For \(t>0\), construct a random finite configuration by first choosing
\(N\) with
\(\mathbb P(N=n)=e^{-t\Lambda}(t\Lambda)^n/n!\), and, conditional on
\(N=n\), taking \(n\) independent marks with law \(\mu/\Lambda\).
The antichain event is unchanged by reordering those marks. Its
conditional probability is \(A_n(\mu)/\Lambda^n\), and hence
\[
 \begin{aligned}
 e^{t\Lambda}\mathbb P(\Pi_t\text{ is an antichain})
 &=e^{t\Lambda}\sum_{n\ge0}
      e^{-t\Lambda}\frac{(t\Lambda)^n}{n!}
      \frac{A_n(\mu)}{\Lambda^n}\\
 &=\sum_{n\ge0}\frac{\mathcal Z_n(\sigma,c)t^n}{n!}
 =G_\mu(t).
 \end{aligned}
\]
This is the finite-intensity Poisson configuration representation of
the source process. At \(t=0\) take the empty configuration almost
surely; both sides are \(1\).

We prove the required spatial independence before using it.
A finite configuration is a finite ordered list; restriction \(R_U\) to a
measurable set \(U\subseteq X\) retains exactly the marks in \(U\),
keeping their original relative order. This restriction is measurable.
Indeed, on the count-\(n\) fiber, split the source into the finitely many
measurable regions
\[
       H_S=\{x\in X^n:x_j\in U\ \Longleftrightarrow\ j\in S\},
                  \qquad S\subseteq\{1,\ldots,n\}.
\]
On \(H_S\), the restriction is the fixed coordinate projection to the
increasing enumeration of \(S\), followed by the measurable inclusion
of the count-\(|S|\) fiber. Preimages are finite unions of these
measurable pieces.

For a finite intensity \(\nu\), define its nonnegative configuration
measure by
\[
 \mathfrak S_{t,\nu}(D)
   =\sum_{n\ge0}\frac{t^n}{n!}\,
     \nu^{\otimes n}(D_n),
 \qquad
 D_n=\{x\in X^n:(n,x)\in D\}.
\]
Its total mass is \(e^{t\nu(X)}\), and the sampled Poisson-count and
independent-mark construction gives the equality of measures
\[
         \mathfrak S_{t,\nu}=e^{t\nu(X)}P_{t,\nu}.
\]
This includes \(\nu=0\): both normalized laws are concentrated on the
empty configuration, while the zero-dimensional product measure has
mass one.

Take any two measurable configuration events \(D,E\). On \(H_S\),
write \(k=|S|\). Splitting the independent mark coordinates into \(S\)
and its complement, in their increasing orders, proves the exact identity
\[
\begin{aligned}
 &\nu^{\otimes n}
   \{x\in H_S:R_U(n,x)\in D,\ R_{U^c}(n,x)\in E\}\\
 &\hspace{20mm}
 = (\nu|_U)^{\otimes k}(D_k)
       (\nu|_{U^c})^{\otimes(n-k)}(E_{n-k}).
\end{aligned}
\]
There is no permutation-invariance assumption on \(D,E\): the increasing
orders fix which coordinates enter each event. The sets \(H_S\) partition
\(X^n\). There are \(\binom nk\) of size \(k\), so summing gives
\[
 \nu^{\otimes n}\{R_U\in D,\ R_{U^c}\in E\}
 =\sum_{k=0}^n\binom nk
       (\nu|_U)^{\otimes k}(D_k)
       (\nu|_{U^c})^{\otimes(n-k)}(E_{n-k}).
\]
Multiply by \(t^n/n!\), use
\(\binom nk/n!=1/(k!(n-k)!)\), and sum nonnegative terms. Reindexing
by \((k,n-k)\) yields
\[
 \mathfrak S_{t,\nu}\{R_U\in D,\ R_{U^c}\in E\}
     =\mathfrak S_{t,\nu|_U}(D)
        \mathfrak S_{t,\nu|_{U^c}}(E).
\]
Since \(\nu(X)=\nu(U)+\nu(U^c)\), the finite positive exponential
normalizers cancel. Therefore
\[
 P_{t,\nu}\{R_U\in D,\ R_{U^c}\in E\}
       =P_{t,\nu|_U}(D)P_{t,\nu|_{U^c}}(E).
\]
Taking one event to be the whole configuration space gives the actual
marginal restriction laws. The displayed rectangle identity is their
joint product law, and thus proves independence.

For any pairwise disjoint measurable sets \(C_i\), this argument gives
independence of all their restrictions, in the precise finite-family
sense. To see this explicitly, restriction of lists is ordinary filtering,
so \(R_U R_V=R_{U\cap V}\). For a chosen \(C_i\), all other
\(C_j\)-restrictions are unchanged by first restricting to \(C_i^c\).
Apply the preceding two-cell identity and induct on the number of sets.
The result is the product of the actual \(P_{t,\nu|_{C_i}}\) laws.
In particular the counts are independent Poisson variables with means
\(t\nu(C_i)\): each restricted process has that count law by its
sampled Poisson-count definition. This remains true for zero-mass cells.

Now let \(F\) be a finite family of disjoint cells, set
\(U=\bigcup_{i\in F}C_i\), and apply this proved independence to the
process of intensity \(t\nu|_U\). A global antichain restricts to an
antichain in every cell. Hence
\[
 P_{t,\nu|_U}(\text{antichain})
     \le\prod_{i\in F}P_{t,\nu|_{C_i}}(\text{antichain}).
\]
Multiplication by \(e^{t\nu(U)}\), and
\(\nu(U)=\sum_{i\in F}\nu(C_i)\), give exactly
\[
             G_{\nu|_U}(t)\le\prod_{i\in F}G_{\nu|_{C_i}}(t).
\]
The preceding binomial calculation iterates to the occupancy factor
\(n!/\prod_i r_i!\) when \(\sum_i r_i=n\); multiplication by
\(t^n/n!\) cancels this factor into the local factorial denominators.
Thus the label-allocation and Poisson-independence descriptions agree
at the level of the actual sampled process.

Index the distinct closed scaled tiles by \(B\).
Assign each tile boundary point to its first containing tile in a
fixed enumeration. This gives measurable disjoint cells \(C_B\)
covering \(E\), with \(C_B\subseteq B\); the difference between \(C_B\)
and \(B\) has floor measure zero. Such boundary assignment is legitimate
because the tiles are countable, their boundaries are null, and the
floor density is absolutely continuous. Set
\(\mu_B=\mu|_{C_B\times\mathbb R}\).
For \(r\ge1\), enlarging \(C_B\) to \(B\) and applying
Lemma~\ref{lem:tile-height} gives
\[
 A_r(\mu_B)
 \le\frac{1}{(2\pi)^{r/2}}
       \int_{\mathcal A_r(B)}e^{-a\|\mathbf h\|^2}
                          \prod_{k=1}^rw(\mathbf f^{(k)})
                          \,d\mathbf f\,d\mathbf h
 \le w_B^r v_r(\varepsilon,c).
\]
The same inequality is equality at \(r=0\). Therefore
\[
 G_{\mu_B}(t)\le
       \sum_{r\ge0}\frac{v_r(\varepsilon,c)(tw_B)^r}{r!}
       =\mathcal V(tw_B).
\]
The finite-family product estimate now implies
\[
 G_{\mu|_{U_m\times\mathbb R}}(t)
       \le\prod_{B\in F_m}\mathcal V(tw_B),
 \qquad U_m=\bigcup_{B\in F_m}C_B,
\]
where \(F_m\) are increasing finite prefixes exhausting all tiles.

We justify both countable limits. First, every finite configuration
has all its floors in some \(U_m\), since every individual floor
belongs to a cell and only finitely many floors occur. For fixed
\(n\), monotone convergence therefore gives
\[
 A_n(\mu|_{U_m\times\mathbb R})\uparrow A_n(\mu).
\]
Applying monotone convergence again to the nonnegative sum over \(n\)
proves
\[
 G_{\mu|_{U_m\times\mathbb R}}(t)\uparrow G_\mu(t).
\]
This is the finite-configuration, or Fock-series, exhaustion used in Lean.

Second, the product on the right is a genuine finite convergent
product. Here are the needed analytic bounds, established independently
of the product inequality. Put \(p=(d-1)/d\in(0,1)\) and \(D_0=2d\).
The proved uniform diameter bound for the actual tiles gives
\(\|\mathbf u-\mathbf v\|\le D_0\) whenever \(\mathbf u,\mathbf v\)
lie in the same tile of scale at most \(1\). Hence
\[
 \|\mathbf v\|^2\ge\tfrac12\|\mathbf u\|^2-D_0^2.
\]
Boundedness of \(\sigma\) thus gives a tile-independent constant \(C\)
such that
\[
       w_B^p\le C e^{-p\|\mathbf u\|^2/4}\qquad(\mathbf u\in B).
\]
Every cell has volume
\(b_\varepsilon=|\varepsilon\Delta|
=\varepsilon^{d-1}|\Delta|>0\). Integrating over the disjoint
cells and using the integrable Gaussian majorant yields
\[
 b_\varepsilon\sum_Bw_B^p
 \le C\int_E e^{-p\|\mathbf u\|^2/4}\,d\mathbf u<\infty.
\]
Lemma~\ref{lem:series}, with its factorial-series proof, supplies an
all-argument exponential upper bound
\(\mathcal V(s)\le A e^{K s^p}\) for \(s\ge0\), where \(A\ge1\)
and \(K>0\) depend on the fixed parameters.
For \(0\le s\le1\), nonnegative coefficients and \(v_0=1\) also give
\[
 1\le\mathcal V(s)\le1+s(\mathcal V(1)-1),\qquad
 \log\mathcal V(s)\le s(\mathcal V(1)-1).
\]
For \(s\ge1\), the exponential bound gives
\(\log\mathcal V(s)\le(\log A+K)s^p\).
Since \(s\le s^p\) for \(0\le s\le1\), these two bounds imply
\[
 0\le\log\mathcal V(s)\le C_0s^p\qquad(s\ge0).
\]
Consequently, for fixed \(t>0\),
\[
 \sum_B\log\mathcal V(tw_B)
       \le C_0t^p\sum_Bw_B^p<\infty.
\]
It follows that the finite products converge to
\[
 \prod_B\mathcal V(tw_B)
   =\exp\!\left(\sum_B\log\mathcal V(tw_B)\right)<\infty.
\]
At \(t=0\) every factor is \(1\), so the same conclusion holds.

Taking the two countable limits in the finite-family bound proves
\[
 \mathcal Z(t)=G_\mu(t)
       =\sum_{n\ge0}\frac{\mathcal Z_n(\sigma,c)t^n}{n!}
       \le\prod_B\mathcal V(tw_B)\qquad(t\ge0).
\]
Finally, when \(t>0\), nonnegativity makes each series term at most
the whole sum:
\[
 \frac{\mathcal Z_n(\sigma,c)t^n}{n!}
       \le\prod_B\mathcal V(tw_B).
\]
Multiplying by \(n!t^{-n}\) gives the claimed occupancy bound,
in particular for every \(n\ge1\).
\end{proof}

\begingroup\small\raggedright
The actual model and Poissonization are checked by
\nolinkurl{Antichain.weightedPointLaw_product},
\nolinkurl{Antichain.weightedPointLaw_antichain},
\nolinkurl{Antichain.weightedPoissonLaw_isProbabilityMeasure}, and
\nolinkurl{Antichain.weightedPoisson_source_expansion}.
The literal filtering maps and fixed-count rectangle expansion are checked by
\nolinkurl{Antichain.measurable_configurationRestriction} and
\nolinkurl{Antichain.pi_configurationRestriction_rectangle}.
The exponential normalization, including zero intensity, is
\nolinkurl{Antichain.configurationSeriesMeasure_eq_exp_smul_poisson_all}.
The genuine restriction law and independence are
\nolinkurl{Antichain.finiteIntensityPoissonLaw_restriction_joint},
\nolinkurl{Antichain.configurationRestriction_composition}, and
\nolinkurl{Antichain.finiteIntensityPoissonLaw_disjoint_restrictions_independent}.
The independent Poisson count laws are
\nolinkurl{Antichain.finiteIntensityPoissonLaw_regionCount} and
\nolinkurl{Antichain.finiteIntensityPoissonLaw_disjoint_counts_joint}.
The finite-union and countable passages use
\nolinkurl{Antichain.configurationSeriesMeasure_restrict_biUnion_le_poisson} and
\nolinkurl{Antichain.configurationSeriesMeasure_countable_le_exp_sum_poisson}.
The one-cell bound is
\nolinkurl{Antichain.actual_tile_configurationPartition_bound}.
Product convergence is checked by
\nolinkurl{Antichain.alcoveWeightSup_rpow_summable},
\nolinkurl{Antichain.tileSeries_log_uniform_bound},
\nolinkurl{Antichain.tileSeries_weighted_logs_summable}, and
\nolinkurl{Antichain.actual_tile_product_hasProd}.
The final statements, with the alcove facts supplied by their direct
proofs and no external mathematical parameter, are
\nolinkurl{Antichain.weightedPartitionSeries_le_tileProduct_direct} and
\nolinkurl{Antichain.weightedPartition_le_factorial_tileProduct_direct}.
\par\endgroup

